\chapter{加入我们的共同体}\label{s:joining}

我们希望你会选择帮我们把这本书做得更好。
如果你对这种协作方式还不熟悉，
请先看\appref{s:conduct}里的行为准则，
然后：

\begin{description}

\item[从小处开始。]
  改一个错字、
  把某道练习的措辞讲清楚、
  更正或更新一条文献引用，
  或者为说明某个观点而建议一个更好的例子或类比。

\item[加入对话。]
  看看别人已经提交的 issue 和改进提案，
  并给它们加评论。
  改进本身往往也可以再改进，
  而且这是向共同体自我介绍、结交新朋友的好办法。

\item[先讨论，再编辑。]
  如果你想提一项大改动，
  比如重组或拆分整整一章，
  请提交一个 issue，概述你的提案和理由，并打上「Proposal」标签。
  我们鼓励每个人都给这些 issue 加评论，
  好让关于「改什么、为什么改」的讨论全部摆在明面上、并可以被存档。
  如果提案被接受，
  实际工作随后可以拆成若干个能独立处理的小 issue 或小改动。

\end{description}

\seclbl{使用这些材料}{s:joining-using}

正如\chapref{s:intro}所说，
本书所有材料都可以在
知识共享 署名—非商业性使用 4.0 许可协议下自由分发和再利用
（\appref{s:license}）。
你可以在任何课堂（免费或收费）中使用 \url{http://teachtogether.tech/} 上的在线版本，
也可以依据\hreffoot{https://en.wikipedia.org/wiki/Fair\_use}{合理使用}条款引用其中的简短片段，
但不能未经事先许可就把大部分内容转载到商业作品中。

这些材料已经被以多种方式使用，
从持续数周的在线课程到密集的线下工作坊。
通常可以在两个漫长的日子里，
覆盖\chapref{s:models}到\chapref{s:process}、
\chapref{s:performance}
和\chapref{s:motivation}的大部分内容。

\subsection*{线下}

这是讲授这项培训最有效的方式，
但要求也最高。
参与者人在一起。
当他们需要以小组练习讲课时，
一部分人或所有人去附近的休息区。
参与者用各自的平板或笔记本电脑查看课上的在线材料、
以及用于共享记笔记（\secref{s:classroom-notetaking}），
其他练习则用纸笔或白板。
提问和讨论都出声进行。

如果你用这种形式授课，
就应该用便利贴当状态旗，
好让你看出谁需要帮助、
谁有问题、
谁准备好往下走（\secref{s:classroom-sticky-notes}）。
你还应该用它们来分配注意力，
让每个人都公平地分到老师的时间；
并把它们当作一分钟卡片，鼓励学习者反思刚学到的东西，
并趁你还有时间行动时，给你可行动的反馈。

\subsection*{分组在线}

在这种形式里，
10—40 名学习者分成 2—6 组、每组 4—12 人，
但这些组在地理上是分散的。
每个组用一台摄像头和麦克风接入视频通话，
而不是每个人都单独在线。
双向通话里，好的音频比好的视频更重要：
只有声音、没有画面（像收音机）比只有画面、没有讲述要好懂得多，
而且只要问题能被清楚地听到，
讲课的人不必能看见人才能回答问题。
话虽如此，
如果一门课不可无障碍使用，它就是没用的（\secref{s:motivation-accessibility}）：
提供描述性文字，在音质差时对每个人都有帮助，
而在音质不差时对有听觉障碍的人也有帮助。

全班一起共享记笔记，
也用共享笔记来提问和答疑。
几十个人同时在通话里说话效果很差，
所以在大多数场次里，
老师负责讲，学习者通过笔记工具的聊天来回应。

\subsection*{个人在线}

分组在线的自然延伸，是个人在线。
和分组在线一样，
老师会讲大部分内容，学习者主要通过文字聊天参与。
好的音频再次比好的视频更重要，
参与者应该用文字聊天来表示自己想下一个发言（\appref{s:meetings}）。

让参与者各自在线，会更难绘制和分享概念图（\secref{s:memory-exercises}）、
或对教学给出反馈（\secref{s:performance-exercises}）。
因此老师应该更多依靠有书面结果、可以放进共享笔记的练习，
比如对现成的教学视频给反馈。

\subsection*{多周在线}

全班每周通过视频会议见一次，每次一小时。
为了照顾学习者的时区和日程，每次聚会可能会办两场。
参与者像上面描述的分组在线课那样使用共享记笔记，
在课间在线提交作业，
并互相评论彼此的作品。
实践中，
评论相对少见：
人们强烈倾向于在每周的聚会上讨论材料。

这是最早使用的一种形式，
我现在不再推荐它：
把课程拉长虽然给了人们时间反思、做更大的练习，
但也大大增加了他们因为其他时间上的要求而不得不中途放弃的几率。

\seclbl{贡献与维护}{s:joining-contributing}

各种贡献都欢迎，
从改进建议到勘误再到新内容。
所有贡献者都必须遵守我们的行为准则（\appref{s:conduct}）；
提交你的作品，
即表示你同意它可以以原始或编辑后的形式被采用，
并以与本书其余材料相同的许可协议发布（\appref{s:license}）。
如果你的材料被采用，
除非你另有要求，我们会把你加到致谢名单里（\secref{s:intro-acknowledgments}）。

本书的源码存放在 GitHub 上：

\begin{center}
  \url{https://github.com/gvwilson/teachtogether.tech/}
\end{center}

\noindent
如果你会使用 Git 和 GitHub，并且想改动、修复或添加点东西，
请提交一个修改 LaTeX 源码的\gref{g:pull-request}{拉取请求}。
如果你想预览自己的改动，
请在命令行里运行 \texttt{make~pdf} 或 \texttt{make~html}。

如果你想报告错误、
提问，
或提出建议，
请在仓库里提交一个 issue。
这么做需要一个 GitHub 账号，
但不需要会使用 Git。

如果你不想创建 GitHub 账号，
请把你的贡献发邮件到 \texttt{gvwilson@third-bit.com}，
并在主题行里带上「T3」或「Teaching Tech Together」。
我们会尽量在一周内回复。

最后，
我们总是乐于听到人们如何使用这些材料，
并且\emph{永远}感激更多的示意图。
